\documentclass[print]{styledarticle}

%%%%%%%%%%%%%%%%%%%%%%%%%%%%%%%%%%%%%
% Load packages here! (or at least, before \ProcessPackages)
\usepackage{extra}
\usepackage{algpseudocode}
\usepackage{algorithm}
\usepackage{multicol}
\usepackage{import}
\usepackage{graphicx}
\usepackage{booktabs}
\usepackage{array}
\usepackage{makecell}
\usepackage{float}
\captionsetup[figure]{hypcap=false}
%%%%%%%%%%%%%%%%%%%%%%%%%%%%%%%%%%%%%

\title{Research Proposal}
\author{Antoni Nikolov \and Catalin Popoiu \and Guilherme Oliveira}
\date{}

\begin{document}

\maketitle

\begin{multicols*}{2}
    \begin{abstract}
        EPF has recently seen a resurgence in interest by companies whose operations are energy-dependent, and therefore benefit from timing their energy use in synchrony with the lowest energy prices. However, the
        literature vastly underreports on the effects of weather data as a key factor for EPF.\\ The use of multi-feature models that include such data also lags behind the current advancements in Transformer and RNN model architectures,
        as well as preprocessing opportunities. This research proposes a plan to test multiple architectures and preprocessing techniques in a modular way.
        We hope to find in its midst, a model that out-performs the current solution used by OpenRemote.
    \end{abstract}

    \import{Sections/}{Introduction}

    \import{Sections/}{LiteratureReview}

    \section{Methodology}
        \import{Sections/}{Data}

        \import{Sections/}{ResearchDesign}

    \import{Sections/}{PreliminaryData}

    \import{Sections/}{StatementOfLimitations}

    \import{Sections/}{HardwareSpecification}

    \import{Sections/}{Conclusion}
\end{multicols*}

\bibliographystyle{plain}
\bibliography{citations}

\import{Sections/}{Appendix}

\end{document}
